\section{Harnais d'expérimentation et bac à sable hermétique}\label{sec:harnais-sandbox}

Cette partie décrit le chantier récent qui a unifié le cycle de vie des harnais
de calibration et ajouté un bac à sable hermétique exécutable hors du
laboratoire. Elle sépare trois objets : l'anatomie d'un harnais, le noyau
partagé qui construit son objet d'exécution, et le bac à sable qui exécute un
bras déterministe sous un garde-fou d'audit. Trois commits en portent la trace
sur la branche \code{dev/1} : \code{11b820c} (bac à sable), \code{02320d5}
(fermeture des trous du garde-fou) et \code{5101523} (noyau partagé).

\subsection{Anatomie d'un harnais}

Un harnais est un jeu de fichiers à responsabilité unique, sous
\code{benchmarks/experiments/}. Le harnais MQTT en donne la forme complète. Le
\textbf{contrat} est \code{mqtt\_reference.py} : il valide le briefing, lie
l'outil à la frontière d'exécution, fournit l'\textbf{adaptateur D1} par le bras
déterministe \code{run\_rules} et recompte le run par \code{evaluate\_trial}. Le
\textbf{bras LLM} est \code{mqtt\_reference\_llm.py} ;
\code{mqtt\_formats\_llm.py} en est la variante H2, qui ne reçoit aucun parseur.
Le \textbf{contrôleur live} est \code{mqtt\_reference\_live.py}, à ne pas
confondre avec le bras : il construit l'agent LLM, l'enveloppe d'un
\code{RecordingClient} qui journalise chaque appel de transport et écrit
\code{decision.json}, \code{cost\_summary.json} et \code{run\_meta.json}. Le
\textbf{déploiement} est \code{mqtt\_reference\_proxmox.py} : la classe
\code{Proxmox} soumet des opérations via \code{ansible-playbook}, la classe
\code{Fixture} démarre le broker depuis l'allocation.

La même anatomie se répète pour SSH, l'API et le contrôle V0, avec une
différence importante : le harnais H2 (\code{mqtt\_formats.py}) ne possède ni
déploiement Proxmox ni client live actif dans le bac à sable, car son corpus est
gelé.

\subsection{Le noyau partagé}

Le noyau existe pour une raison précise, consignée dans son en-tête : la
frontière de gouvernance de \code{src/agent/core/executor.py} lit un ensemble
fixe d'attributs par \code{getattr} avec valeur par défaut. Une omission ne se
manifeste donc qu'au moment où un appel d'outil se comporte mal. Chaque harnais
reconstruisait cet objet à la main ; le noyau remplace les copies par un unique
\code{build\_run}.

\code{REQUIRED\_ATTRIBUTES} énumère treize attributs : le dossier du run, les
trois budgets, le verrou d'archivage, le verrou de preuve, le compteur d'appels,
l'horloge de départ, le motif de limite, puis la portée, le contexte, la
politique et le suivi de coût. Un attribut mérite un commentaire particulier :
\code{\_evidence\_integrity\_failed}. \code{check\_execution\_limits} le lit en
première instruction et lève \code{EvidenceWriteError} s'il est vrai ;
\code{wrap\_tool.archive} le pose quand l'archive ne peut pas être écrite. Ce
verrou documente une règle : une preuve perdue doit rester visible, et le run
devient non vérifiable. Un harnais qui l'oublierait retomberait sur la valeur
par défaut et poursuivrait un run dont l'archive est trouée.

Le refus est à deux niveaux : \code{build\_run} vérifie l'attribut manquant et
la clé de contexte absente avant de rendre l'objet, et
\code{require\_attributes} rejoue le même contrôle sur un objet quelconque. Une
variante déterministe ne déclare ni modèle ni fournisseur.

La couverture reste incomplète : les quatre harnais \code{*\_live.py} passent
par une fonction locale homonyme \code{new\_run}, et \code{new\_run} du noyau
n'a aucun appelant direct ; la délégation passe par \code{build\_run}.

\subsection{Le bac à sable}

\code{benchmarks/experiments/sandbox.py} prouve, sur un poste de développement
comme en intégration continue, qu'un harnais s'exécute sans broker, sans modèle
et sans laboratoire. Il charge le corpus gelé (huit fixtures) et, pour chaque
dossier : valide le dossier, ouvre un journal, exécute le bras choisi, écrit les
artefacts, puis recompte par l'évaluateur existant. Cinq artefacts sont produits
par essai : \code{dossier.json}, \code{oracle.json}, \code{answer.json},
\code{journal.jsonl} et \code{evaluation.json}. Un
\code{sandbox-report.json} agrège les essais et le bloc \code{guard}.

Le bras est un paramètre explicite, et un nom inconnu lève une erreur. Ce choix
est une correction : la version d'origine lisait le bras depuis le corpus, qui
ne le déclare pas, rendant le bras par convention inatteignable et faisant
inventer au rapport la politique testée.

\subsection{Le garde-fou d'audit et ses limites}

Le garde-fou repose sur \code{sys.addaudithook}. La classe \code{Guard} reste
installée dans le processus, mais n'inspecte les événements que lorsqu'elle est
armée : le gestionnaire \code{active()} borne l'application au run.
\code{install\_guard} n'installe le hook qu'une fois et rebinde ensuite le
dossier autorisé et la liste des violations à chaque run ; sans ce rebind, un
second run observerait le dossier du premier et étoufferait ses propres
violations. Les événements bloqués couvrent les sockets, les sous-processus, la
famille \code{exec}/\code{spawn}/\code{posix\_spawn}, \code{ctypes} et les
opérations de renommage et de suppression. Les écritures sont traitées à part :
l'événement \code{open} est refusé si les drapeaux contiennent un mode
d'écriture et si le chemin résolu sort du dossier autorisé.

L'honnêteté de ce mécanisme se lit au commit \code{02320d5}. Un audit en lecture
seule a cassé la première version par deux échappées mesurées :
\code{os.posix\_spawn} démarrait un processus, et \code{ctypes.CDLL(None).system()}
exécutait du code natif, alors que le rapport annonçait une exécution hermétique
sans violation. Le rapport a donc été corrigé sur trois points : la clé
\code{hermetic} est remplacée par \code{guard\_clean}, qui ne dit que l'absence
d'événement listé ; un bloc \code{guard} déclare le mécanisme, les événements
bloqués et, explicitement, les angles morts ; et \code{claims\_isolation} est
fixé à \code{false}. Ce dernier indicateur est faux par construction : un hook
d'audit observe des événements de l'interpréteur Python, pas des appels système.
Il ne voit pas le code natif atteint hors \code{ctypes}, ni une écriture par un
descripteur ouvert avant l'armement, ni un appel système émis hors CPython.

Cette limite a une conséquence méthodologique : \code{guard\_clean: true}
signifie « aucun des événements listés n'est survenu », jamais « le run est
isolé ».

\begin{table}[H]
\centering
\caption{Ce que le garde-fou couvre, et ce qu'il laisse hors d'observation.}
\label{tab:guard-limits}
\small
\begin{tabular}{@{}p{6.2cm}p{7.6cm}@{}}
\toprule
Couvert, avec refus levé & Hors d'observation \\
\midrule
Sockets : connexion, résolution, envoi, écoute & Code natif atteint sans \code{ctypes}, \code{Popen} ou \code{posix\_spawn} \\
Processus : \code{subprocess}, \code{os.system}, \code{exec*}, \code{fork}, \code{spawn*}, \code{posix\_spawn*} & Écritures via un descripteur ouvert avant l'armement \\
\code{ctypes} : chargement, résolution de symbole, appel & Appels système émis hors de l'interpréteur CPython \\
Renommage et suppression : \code{os.rename}, \code{remove}, \code{replace}, \code{unlink}, \code{shutil} & Tout comportement après désarmement : le hook reste installé mais inerte \\
Écritures par \code{open} hors du dossier autorisé (chemins résolus) & Aucune garantie d'isolation du processus \\
\bottomrule
\end{tabular}
\end{table}

\subsection{Ce que le bac à sable démontre}

Sur les huit fixtures gelées, le bras à parser déclaré réussit six essais sur
huit : les deux échecs sont les essais \code{variant/raw}, où le parseur
déclaré ne couvre pas le schéma et s'abstient sur chaque message, d'où une
conclusion fautive. Le bras par convention, qui lit les clés sans schéma
déclaré, réussit huit sur huit. La réussite n'est pas comptée sur la conclusion
seule : l'évaluateur l'impose comme conclusion correcte, sans contradiction et
sans aucun verdict erroné.

Les deux bras produisent la même leçon hors ligne que H2 : l'échec sur
\code{variant/raw} ne vient pas d'un manque d'intelligence mais d'une frontière
du parseur déclaré, et un script qui essaie rétablit la couverture. Le bac à
sable instrumente le contraste ; il ne le mesure pas sur un modèle. Aucun bras
LLM n'y figure, aucune attestation de pré-vol n'est produite, et ni
\code{pair\_admissible} ni \code{mt\_admissible} n'en sortent.

\subsection{Limites du bac à sable}

Cinq limites doivent être lues avec les résultats précédents. D'abord, le bac à
sable ne met en jeu aucun modèle : sa réussite est une propriété du harnais, pas
un résultat expérimental, et le rapport porte \code{claims\_measurement: false}.
Ensuite, \code{guard\_clean} est une observation de Python et non une preuve
d'isolation. Troisièmement, la portée est partielle : les tests négatifs qui
prouvent que le hook refuse réellement réseau, processus, \code{posix\_spawn} et
\code{ctypes} passent par l'appel direct du garde, et aucun test de bout en bout
ne confirme que le run déclenche effectivement le hook installé ; le chemin de
détection n'est prouvé qu'au niveau unitaire. Quatrièmement, le corpus est gelé,
n'a qu'une scénario et huit fixtures, ce qui suffit à instrumenter le contraste
mais pas à estimer un taux. Enfin, le bac à sable contourne la frontière de
gouvernance qu'il est censé accompagner : il appelle le bras directement, sans
construire de run ni passer par l'enveloppe d'outils, et le journal est écrit
par le harnais lui-même ; le noyau partagé et le garde-fou restent deux
instruments séparés, chacun vérifié isolément.

\subsection{Tests}

L'ensemble compte vingt-quatre tests : quatorze pour le bac à sable et dix pour
le noyau. Ils prouvent quatre propriétés. Pour le noyau, le run expose la
surface de gouvernance, un run incomplet est refusé à la construction,
\code{check\_execution\_limits} accepte un run neuf, et le dépassement du budget
d'appels est levé par l'enveloppe d'outils et non par le contrôle de limites.
Pour le bac à sable, les huit fixtures produisent les cinq artefacts et un
rapport cohérent, et le bras est déclaré et jamais lu depuis le corpus. Pour la
leçon H2, le parser déclaré échoue sur au moins un essai quand le bras par
convention réussit partout. Pour le garde-fou, des tests négatifs refusent les
écritures hors dossier, le réseau, les sous-processus, \code{posix\_spawn} et
\code{ctypes}, tandis qu'un test vérifie que le hook reste inerte après
désarmement et qu'un second garde observe son propre dossier.

Ce que ces tests ne prouvent pas : aucun n'exécute un bras LLM, aucun ne mesure
d'intervalle de confiance, et aucun ne vérifie que le hook installé est atteint
pendant le run plutôt que seulement quand il est appelé à la main.
